\documentclass[10pt]{article}
\providecommand{\FIGDIR}{fig}
\newcommand{\DIGESTDATE}{7 August 2026}
% ===================== PM Research Watch - shared preamble =====================
\usepackage[T1]{fontenc}
\usepackage[utf8]{inputenc}
\usepackage{amsmath,amssymb}
\usepackage{tgpagella}
\usepackage{tgheros}
\renewcommand{\sfdefault}{qhv}
\renewcommand{\ttdefault}{lmtt}
\usepackage{microtype}
\usepackage[english]{babel}

\usepackage[a4paper,top=17mm,bottom=17mm,left=17mm,right=17mm,
            headsep=6mm,footskip=10mm]{geometry}
\usepackage{graphicx}
\usepackage{xcolor}
\usepackage{colortbl}
\usepackage{tikz}
\usetikzlibrary{calc}
\usepackage[most]{tcolorbox}
\usepackage{titlesec}
\usepackage{fancyhdr}
\usepackage{booktabs}
\usepackage{array}
\usepackage{longtable}
\usepackage{enumitem}
\usepackage{multicol}
\usepackage{ragged2e}
\usepackage[hidelinks]{hyperref}

% ------------------------------------------------------------------ palette
\definecolor{Ink}    {HTML}{14262E}
\definecolor{Deep}   {HTML}{0F4C5C}
\definecolor{Teal}   {HTML}{1F7A8C}
\definecolor{Sky}    {HTML}{4FA3C4}
\definecolor{Sage}   {HTML}{6FA287}
\definecolor{Amber}  {HTML}{E3A008}
\definecolor{Clay}   {HTML}{B5643C}
\definecolor{Coral}  {HTML}{D1495B}
\definecolor{Violet} {HTML}{6A5B9E}
\definecolor{Slate}  {HTML}{4A6572}
\definecolor{Paper}  {HTML}{FBF9F5}
\definecolor{Mist}   {HTML}{EEEAE2}
\definecolor{Rule}   {HTML}{D9D3C9}

\pagecolor{Paper}
\color{Ink}

\hypersetup{colorlinks=true, linkcolor=Deep, urlcolor=Teal, citecolor=Teal,
            pdftitle={PM Research Watch}, pdfauthor={Automated daily digest}}

% ------------------------------------------------------------------ headings
\titleformat{\section}
  {\sffamily\bfseries\large\color{Deep}}{}{0pt}
  {}[{\vspace{1.2mm}\color{Amber}\titlerule[1.1pt]}]
\titlespacing*{\section}{0pt}{5.5mm}{2.8mm}

\titleformat{\subsection}
  {\sffamily\bfseries\normalsize\color{Teal}}{}{0pt}{}
\titlespacing*{\subsection}{0pt}{3.6mm}{1.4mm}

\setlist[itemize]{leftmargin=4.2mm, itemsep=1.1mm, topsep=1.3mm, parsep=0pt,
                  label=\textcolor{Teal}{\footnotesize$\blacktriangleright$}}

% ------------------------------------------------------------------ header/footer
\pagestyle{fancy}
\fancyhf{}
\renewcommand{\headrulewidth}{0pt}
\renewcommand{\footrulewidth}{0pt}
\fancyhead[L]{\sffamily\scriptsize\color{Slate}PM RESEARCH WATCH\;\textcolor{Rule}{\textbar}\;\DIGESTDATE}
\fancyhead[R]{\sffamily\scriptsize\color{Slate}Particulate matter monitoring \& health effects}
\fancyfoot[C]{\sffamily\scriptsize\color{Slate}\thepage}

% ------------------------------------------------------------------ components
% metric tile
\newcommand{\metric}[3]{%
  \begin{tcolorbox}[enhanced, colback=white, colframe=Rule, boxrule=0.5pt,
    arc=1.6mm, left=2.2mm, right=2.2mm, top=1.8mm, bottom=1.8mm,
    borderline west={1.7mm}{0pt}{#3}, nobeforeafter, valign=center,
    height=17mm, width=\linewidth]
    \RaggedRight
    {\sffamily\bfseries\fontsize{17}{17}\selectfont\color{#3}#1}\\[0.5mm]
    {\sffamily\fontsize{6.5}{7.4}\selectfont\color{Slate}#2}
  \end{tcolorbox}}

% pull-quote / key insight box
\newtcolorbox{keybox}[1][]{enhanced, colback=Mist, colframe=Mist,
  arc=1.6mm, left=3.2mm, right=3.2mm, top=2.4mm, bottom=2.4mm,
  borderline west={1.7mm}{0pt}{Amber}, #1}

% subtopic band
\newcommand{\band}[2]{%
  \vspace{3.2mm}\noindent
  \begin{tcolorbox}[enhanced, colback=#2, colframe=#2, arc=1.2mm,
    left=2.6mm, right=2.6mm, top=1.1mm, bottom=1.1mm, boxrule=0pt,
    width=\linewidth, nobeforeafter]
    {\sffamily\bfseries\small\color{white}#1}
  \end{tcolorbox}\vspace{1.6mm}}

% one paper entry: {short cite}{title}{journal}{doi}{body}
\newcommand{\paperentry}[5]{%
  \par\vspace{1.9mm}\noindent
  {\sffamily\bfseries\footnotesize\color{Ink}#1}\;%
  {\sffamily\fontsize{7}{8}\selectfont\color{Slate}#3}\\[0.4mm]
  {\footnotesize\itshape\color{Deep}#2}\\[0.6mm]
  {\footnotesize #4}\\[0.3mm]
  {\fontsize{6.6}{7.6}\selectfont\color{Slate}\texttt{doi:}\,\href{https://doi.org/#5}{\texttt{\detokenize{#5}}}}
  \par}

% caption style
\newcommand{\figcap}[1]{\par\vspace{1.0mm}
  {\fontsize{7.2}{8.6}\selectfont\color{Slate}\textsf{\textbf{Fig.}}\; #1\par}\vspace{1.5mm}}

\providecommand{\permil}{\textperthousand}

\setlength{\parindent}{0pt}
\setlength{\parskip}{0pt}
\renewcommand{\arraystretch}{1.18}


\begin{document}

% ==================================================================== MASTHEAD
\thispagestyle{empty}
\begin{tikzpicture}[remember picture, overlay]
  \fill[Deep] (current page.north west) rectangle
      ($(current page.north east)+(0,-67mm)$);
  \fill[Amber] ($(current page.north west)+(0,-67mm)$) rectangle
      ($(current page.north east)+(0,-68.4mm)$);
  \foreach \x/\y/\r in {22/12/0.9, 41/26/1.5, 63/9/0.7, 88/31/1.1, 112/17/1.9,
                        138/29/0.8, 159/13/1.3, 178/54/1.0, 196/21/1.6,
                        30/44/0.6, 74/57/1.2, 125/48/0.9, 168/6/0.8, 191/59/0.7,
                        52/18/1.0, 100/60/0.6, 148/20/1.1, 66/49/1.4, 205/45/0.8}
    \fill[Sky, opacity=0.30] ($(current page.north west)+(\x mm,-\y mm)$) circle (\r mm);
\end{tikzpicture}

\vspace*{4mm}
{\sffamily\fontsize{7.6}{9}\selectfont\color{Sky}%
 AUTOMATED DAILY LITERATURE DIGEST\;\textbullet\;PhD RESEARCH WATCH\par}
\vspace{2.2mm}
{\sffamily\bfseries\fontsize{27}{29}\selectfont\color{white}%
 Particulate Matter\par}
{\sffamily\fontsize{27}{29}\selectfont\color{Sky}%
 Monitoring \& Health Effects\par}
\vspace{5.0mm}
{\sffamily\fontsize{8.4}{10}\selectfont\color{white}%
 Issue 2026-08-07\;\;\textcolor{Amber}{\textbar}\;\;Entry window 7 August 2026\;\;%
 \textcolor{Amber}{\textbar}\;\;26 records in scope\par}
\vspace{18mm}

% -------------------------------------------------------------- metric strip
\noindent
\begin{minipage}[t]{0.190\linewidth}\metric{26}{RECORDS IN SCOPE\\(20 rejected, 4 metadata-only)}{Deep}\end{minipage}\hfill
\begin{minipage}[t]{0.190\linewidth}\metric{7}{SENSING RECORDS\\(27\% OF CORPUS)}{Teal}\end{minipage}\hfill
\begin{minipage}[t]{0.190\linewidth}\metric{8}{SUBTOPIC CLUSTERS\\REPRESENTED}{Sky}\end{minipage}\hfill
\begin{minipage}[t]{0.190\linewidth}\metric{8}{TIER-A\\RECORDS}{Amber}\end{minipage}\hfill
\begin{minipage}[t]{0.190\linewidth}\metric{16}{EFFECT ESTIMATES\\EXTRACTED}{Clay}\end{minipage}

\vspace{6mm}

% -------------------------------------------------------------- signal of day
\section*{\sffamily\bfseries\color{Deep}Signal of the day}
\vspace{-1.5mm}{\color{Amber}\rule{\linewidth}{1.1pt}}\vspace{2.5mm}

\begin{keybox}
{\sffamily\bfseries\small\color{Deep}Two \textit{EHP} cohorts published on the same day
both find that \emph{where the particle came from} carries more information than how much
of it there is --- and one of them shows the effect survives inside the household, where
no monitoring network reaches.}\\[1.4mm]
\footnotesize
Cantuaria et al.\ pooled \textbf{159{,}769} participants from the Danish Diet, Cancer and
Health cohort and the Danish National Health Survey, reconstructed 10-year address
histories, and decomposed exposure into \emph{local traffic} and \emph{other} sources for
each of PM$_{2.5}$, NO$_2$, UFP and elemental carbon. Total PM$_{2.5}$ gave the largest
association with incident COPD (\textbf{HR 1.11, 95\% CI 1.05--1.17} per IQR), and the
source split inverts between metrics: PM$_{2.5}$ \emph{from sources other than local
traffic} was the more harmful fraction, whereas for UFP and EC the local-traffic
contribution dominated. A network built to resolve traffic gradients is therefore
optimised for the wrong fraction of the PM$_{2.5}$ signal and the right one for UFP.
Sassano et al.\ then supply the indoor mirror image in the Golestan Cohort
(\textbf{50{,}045} adults, 724{,}064 person-years, 262 kidney deaths): biomass household
fuel carried \textbf{HR 1.20 (1.04--1.37)} per 10 years of use for kidney-disease
mortality, kerosene and gas were null --- and the biomass estimate splits cleanly on
ventilation, \textbf{1.19 (1.06--1.34)} without a chimney against \textbf{1.06
(0.95--1.18)} with one ($P_{\text{diff}}=0.025$). The exposure that matters is not
measured by any ambient monitor and is modified by a variable no exposure model carries.
Both results point the same way: source and microenvironment, not concentration, are the
unmeasured axes.
\end{keybox}

\vspace{3mm}

% Bottom margin proofed: at 2 baselines the last bullet closes clear of the folio.
% 2026-08-04 log records that 4-7 baselines start printing over the page number, and
% 2026-08-06 records that the macro must sit BEFORE the section heading to take effect.
\enlargethispage{3\baselineskip}
% -------------------------------------------------------------- exec brief
\section*{\sffamily\bfseries\color{Deep}What changed today}
\vspace{-1.5mm}{\color{Amber}\rule{\linewidth}{1.1pt}}\vspace{2.5mm}

\footnotesize
\begin{itemize}
\item \textbf{Sixteen effect estimates with intervals --- the largest set this watch has
carried}, against six on 6 August and one on 5 August. Fourteen are PM or household-smoke
estimates and two are co-pollutant. The CI drought that ran 2--5 August is now
comprehensively over, and the reason is compositional: three large cohort or time-series
papers (Cantuaria, Orr, Sassano) landed in the same window, each reporting a family of
stratified estimates rather than a single headline. Three GBD burden records landed with
them --- India stillbirths, LRI and TB, ischaemic stroke at 20--54 --- and contribute none,
attributing fractions to air pollution as a risk factor rather than estimating an
exposure--response.

\item \textbf{A well-powered internal contradiction on wildfire smoke and influenza.}
Orr et al.\ (\textit{EHP}) find wildfire-season PM$_{2.5}$ associated with higher
influenza risk in the four states reporting \emph{laboratory-confirmed influenza} ---
Colorado \textbf{RR 1.067 (1.056--1.078)}, Arizona 1.061 (1.026--1.100), Oregon 1.049
(1.041--1.057), Montana 1.038 (1.013--1.063) per 10\,$\mu$g/m$^3$ --- and \emph{no} effect
in Nevada (1.005, 0.920--1.097) and a \textbf{protective} estimate in Washington
(\textbf{0.884, 0.842--0.919}), the two states contributing only syndromic ILI counts.
The split falls exactly along the outcome-ascertainment boundary, which is the most
parsimonious explanation and one the paper raises but does not resolve. Read as an
outcome-misclassification result rather than a geography result.

\item \textbf{PM$_{2.5}$ composition beats PM$_{2.5}$ mass again, this time with a
genotype interaction.} Li et al.\ (\textit{Environ Int}) report that an IQR increase in
all 15 measured PM$_{2.5}$ components jointly was associated with \textbf{10.0\% higher
MRI hepatic fat fraction} in 113 Los Angeles Latino adolescents with obesity, with
elemental carbon, bromine, copper, potassium and lead as the drivers --- a
combustion/traffic signature --- and effects concentrated in \textbf{PNPLA3 GG carriers},
who also showed 550 mediating metabolite features against 196 in non-carriers. Total
mediation was significant \emph{only} in GG carriers. Replication in 81 mixed-ethnicity
young adults held for EC and Br. Small discovery cohort, single metropolitan area, and
component estimates from a regional model rather than personal measurement.

\item \textbf{Instrumentation: the day's best sensing content is largely un-summarisable.}
Four Crossref by-ISSN deposits directly on the sensing axis --- a low-cost PM sensor
performance evaluation under Northeast China pollution episodes, an explainable-ML
attribution of winter PM$_{2.5}$ drivers in China, surfactant effects on supermicron NaCl
hygroscopic growth, and metal-ion/pH control of brown-carbon absorption --- \textbf{carry
no abstract in Crossref, Europe PMC or Semantic Scholar}, so they are listed on metadata
and not summarised. What \emph{is} summarisable is Jin et al.\ (\textit{ACP}):
three-or-more-day-old smoke transported into Missoula raised CO, PM$_{2.5}$ and total VOCs
by factors of \textbf{2--8}, hourly PM$_{2.5}$ peaking at \textbf{120\,$\mu$g/m$^3$},
tested against GEOS-Chem.

\item \textbf{Two low-cost-sensor network records enter as backfill, both about deployment
rather than accuracy.} DeMarsh et al.\ show a community PurpleAir network in Fresno County
flagging four weeks of NAAQS-exceeding PM$_{2.5}$ at locations the regulatory network
missed, while agreeing on concentration where the two overlap. Blanco-Villafuerte et al.\
describe a \textbf{176-monitor} national PurpleAir network in Peru feeding a hybrid
LCS/satellite/CTM model that reached \textbf{R$^2$ 0.88} in Lima. Both are Crossref-dated
January--February 2026: backfill, off any trend line.

\end{itemize}

\clearpage

% ==================================================================== ANALYTICS
\section*{\sffamily\bfseries\color{Deep}Corpus analytics}
\vspace{-1.5mm}{\color{Amber}\rule{\linewidth}{1.1pt}}\vspace{3mm}

\noindent\includegraphics[width=\linewidth]{\FIGDIR/f1_subtopics.png}
\figcap{Subtopic assignment of the 26 in-scope records, single-label by dominant
contribution. \textit{Exposure assessment \& modelling} and \textit{Sensing, forecasting
\& instrumentation} tie at 7 each (27\% apiece), so \textbf{54\% of the window is
methodological rather than epidemiological}. Respiratory \& allergic follows at 4 and
burden/policy at 3. Eight of the ten standing clusters are populated; \textit{Neuro /
mental health} and \textit{Occupational \& indoor} are empty as dominant labels, though
the Sassano household-fuel cohort and the Hou filtration paper both sit adjacent to the
latter.}

\vspace{2mm}
\noindent\includegraphics[width=\linewidth]{\FIGDIR/f2_design_endpoint.png}
\figcap{Left --- study architecture. Modelling/inventory and measurement campaigns tie at
7 each, ahead of chamber/laboratory (3) and acute observational designs (3); only 2
records are cohort studies, but those 2 (Cantuaria, Sassano) carry 210{,}000 participants
between them and supply half the day's effect estimates. Right --- health endpoint by
organ system. \textbf{Fourteen of the 26 records carry no human health endpoint}, and the
twelve that do are spread one per endpoint with no two records sharing a label ---
consistent with a methods-heavy window and with the register's breadth. There is no trial
or intervention record in this window.}

\vspace{2mm}
\noindent\includegraphics[width=\linewidth]{\FIGDIR/f3_metric_geography.png}
\figcap{Left --- particle metric. Eight records use PM$_{2.5}$ alone, two pair it with
PM$_{10}$ and one uses PM$_{10}$ alone. The 15-record unspeciated/emissions bar is genuine
heterogeneity, not missing detail: it holds PM$_{0.3}$ filtration efficiency, speciated
organosulfates, brown-carbon absorption, hygroscopic growth factors, trace-metal
bioaccumulation, household fuel type and three GBD risk-factor attributions. Right ---
geographic provenance. North America leads (6), then the global/multi-region group (7,
which includes the four laboratory and computational records placed there by intent) and
China (4). \textbf{South Asia contributes 3 and Sub-Saharan Africa 1.}}

\vspace{2mm}
\noindent\includegraphics[width=\linewidth]{\FIGDIR/f6_lifecourse.png}
\figcap{Life-course windows addressed; counts are non-exclusive. Ten of 26 records carry a
human life-course window and adults dominate (6). \textbf{The adolescent window is
populated for the first time in five issues} by Li et al.'s PNPLA3 hepatic-fat cohort. The
\textbf{older-adult window is empty for a fourth consecutive issue}, which remains the
most conspicuous standing gap given that group's cardiovascular attributable burden.}

\vspace{2mm}
\noindent\includegraphics[width=\linewidth]{\FIGDIR/f4_heatmap.png}
\figcap{Subtopic $\times$ study-architecture cross-tabulation. The \textit{Sensing} row is
again confined to measurement campaigns and laboratory work with \textbf{no cell in any
observational-epidemiology column}, unbroken in every issue since 27 July. The converse
also holds: the two cohort studies sit entirely in the respiratory and other-clinical
rows. The one crossing cell is Orr et al.'s in-home monitor pilot, which is a sensing
record with a clinical population attached.}

\vspace{2mm}
\noindent\includegraphics[width=\linewidth]{\FIGDIR/f5_forest.png}
\figcap{All sixteen effect estimates carrying an interval, on a common logarithmic axis.
\textbf{Exposure contrasts are heterogeneous and the estimates are not mutually
comparable}: Cantuaria is per IQR of a 10-year mean, Orr is per 10\,$\mu$g/m$^3$ of
wildfire-season PM$_{2.5}$, Sassano is per 10 years of household fuel use, and Hussein is
a threshold contrast against the Malaysian standard. \textbf{The Washington ILI estimate
(0.884, 0.842--0.919) is the only point below the null} and is best read as
outcome-misclassification rather than protection --- it and Nevada are the two
syndromic-ILI states. One entry is a co-pollutant estimate (NO$_2$, Cantuaria) and is
marked as such; the remaining fifteen are PM or household-smoke exposures.}

\clearpage

% ==================================================================== DIGEST
\section*{\sffamily\bfseries\color{Deep}Paper-level digest}
\vspace{-1.5mm}{\color{Amber}\rule{\linewidth}{1.1pt}}\vspace{1mm}

{\fontsize{7.4}{9}\selectfont\color{Slate}Ordered by cluster. Each entry gives the
methodological core and the finding that matters, not an abstract paraphrase.\par}

\band{Respiratory \& allergic\hfill 4 records}{Coral}

\paperentry{Cantuaria et al.\ 2026}
{Source-specific PM$_{2.5}$, UFP, EC and NO$_2$ and incident COPD in two pooled Danish cohorts}
{Environ Health Perspect}
{159{,}769 participants from the Diet, Cancer and Health cohort and the Danish National
Health Survey, with complete residential address histories driving 10-year time-weighted
exposures split into local-traffic and other-source contributions; Cox models adjusted for
demographic, socioeconomic and lifestyle factors including smoking. Per IQR of total
exposure: \textbf{PM$_{2.5}$ HR 1.11 (1.05--1.17)}, NO$_2$ 1.08 (1.04--1.13), UFP 1.05
(0.99--1.11), EC 1.02 (1.00--1.05). \textbf{The source split inverts across metrics} ---
non-traffic PM$_{2.5}$ was the more harmful fraction while traffic-derived UFP and EC
dominated theirs. Associations were stronger in women, never-smokers and prior asthma.
Residual confounding by smoking intensity is the live threat given COPD's dominant cause.
Directly relevant to siting: a traffic-gradient network resolves the wrong PM$_{2.5}$
fraction and the right UFP one.}
{10.1021/EHP.6c00128}

\paperentry{Orr et al.\ 2026 (EHP)}
{Wildfire-season PM$_{2.5}$ and influenza / ILI risk across six western US states}
{Environ Health Perspect}
{Generalised linear distributed-lag models on state health-department influenza and ILI
counts from Arizona, Colorado, Montana, Nevada, Oregon and Washington, 2010--2019, testing
both preceding-wildfire-season exposure on the following flu season and short-term
exposure within it. Per 10\,$\mu$g/m$^3$: \textbf{CO 1.067 (1.056--1.078), AZ 1.061
(1.026--1.100), OR 1.049 (1.041--1.057), MT 1.038 (1.013--1.063)} --- all four states with
laboratory-confirmed influenza data --- against \textbf{NV 1.005 (0.920--1.097) and WA
0.884 (0.842--0.919)}, the two with syndromic ILI only. \textbf{The effect boundary
coincides exactly with the outcome-ascertainment boundary}, so the geographic reading is
almost certainly the wrong one. Ecological design, no individual exposure, and
wildfire-attributable PM$_{2.5}$ is itself a modelled quantity.}
{10.1021/EHP.6c00326}

\paperentry{Hussein et al.\ 2026}
{PM$_{10}$ and emergency asthma presentations in Petaling Jaya, Malaysia}
{J Glob Health}
{Four years of ED records (Jan 2014 -- Apr 2017) from a Klang Valley hospital against
pollutant and meteorological data from nearby stations; generalised additive models plus
distributed lag non-linear models for delayed and non-linear PM$_{10}$ effects. Asthma
exacerbations were 2.3\% of all ED visits. Relative to the Malaysian standard, daily mean
PM$_{10}$ of \textbf{200\,$\mu$g/m$^3$ was associated with a 21.3\% increase
(95\% CI 2.2--44.0)} and \textbf{300\,$\mu$g/m$^3$ with 34.0\% (4.7--71.7)}; the cumulative
lag 0--7 exposure--response turns upward from about 150\,$\mu$g/m$^3$, and the lag
structure shifts from days 2--3 at 100\,$\mu$g/m$^3$ to same-day at $\geq$200. Single
hospital, single monitoring catchment, wide intervals --- the authors correctly label the
associations exploratory. The threshold shape is the transferable part: it argues for
alert triggers well below current national standards.}
{10.7189/jogh.16.04249}

\paperentry{Jiang et al.\ 2026}
{Short-term PM and respiratory-infection hospitalisation burden in kidney transplant recipients}
{J Cent South Univ Med Sci}
{Retrospective time-series over 874 kidney transplant recipients hospitalised for
respiratory tract infection at the Third Xiangya Hospital, 2017--2023, with high-resolution
daily exposure to six pollutants over the 7 days before symptom onset and DLNMs for
lagged and cumulative associations. \textbf{PM$_{2.5}$ showed the strongest association
with burden: at 160\,$\mu$g/m$^3$ cumulative exposure, hospitalisation cost rose by
9{,}544.5 yuan (95\% CI 2{,}548.8--21{,}637.7)}; PM$_{10}$ showed a dose-dependent trend on
length of stay (RR 1.63 at 300\,$\mu$g/m$^3$) whose interval, 0.18--15.94, is so wide it
carries no information and is deliberately excluded from the forest plot. SO$_2$ above
75\,$\mu$g/m$^3$ also raised burden; CO, NO$_2$ and O$_3$ were null. Single centre,
immunosuppressed population, ambient exposure assigned without residential history.
Economic rather than clinical endpoints, which is unusual and useful.}
{10.11817/j.issn.1672-7347.2026.260089}

\band{Sensing, forecasting \& instrumentation\hfill 7 records}{Amber}

\paperentry{Jin et al.\ 2026}
{Emissions, chemistry and health impact of aged wildfire smoke over Missoula, Montana}
{Atmos Chem Phys}
{Hourly surface observations of PM$_{2.5}$, CO, NO$_x$, O$_3$ and \textbf{75 speciated
VOCs} through a strong 2020 smoke episode, used as a direct test of how wildfire emissions,
ageing chemistry and health effects are represented in GEOS-Chem. Smoke three or more days
old, transported from California and the Pacific Northwest, raised CO, PM$_{2.5}$ and total
measured VOCs by \textbf{factors of 2--8}, with hourly maxima of 800\,ppb CO and
\textbf{120\,$\mu$g/m$^3$ PM$_{2.5}$}. The value here is the speciated VOC suite alongside
mass: it is the composition record against which a mass-only low-cost network's health
interpretation has to be judged, and it is one city and one episode, so the transport
chemistry is not generalisable without replication at other receptor sites.}
{10.5194/acp-26-11047-2026}

\paperentry{Orr et al.\ 2026 (JMIR)}
{In-home digital air quality monitors in families of children with asthma: feasibility and acceptability}
{JMIR Pediatr Parent}
{Pilot mixed-methods study, \textbf{N = 20} families, combining interviews and surveys to
assess whether a consumer in-home air quality monitor can be implemented in paediatric
asthma care. Use was reported as both feasible and acceptable, with \textbf{affordability
identified as the principal barrier to clinical adoption}. The report is very short and
carries no exposure data, no adherence metric and no asthma outcome, so it establishes
only that the deployment step is not the bottleneck. That is still the piece of the
low-cost-sensing translation chain that is least often evidenced --- most LCS literature
stops at accuracy and never asks whether a clinical population will host the device. Needs
a powered trial with objective device-uptime data before it can bear weight.}
{10.2196/92070}

\paperentry{Adeniji et al.\ 2026}
{Twelve years of hourly criteria-pollutant and meteorological coupling in Abeokuta, Nigeria}
{PLOS Glob Public Health}
{Hourly PM$_{2.5}$, PM$_{10}$, NO$_2$, SO$_2$, CO and O$_3$ over 2012--2023 with six
meteorological covariates. \textbf{PM$_{2.5}$ and PM$_{10}$ dominate and persistently
exceed WHO guidelines}, peaking in the December--March Harmattan; regression against
temperature gives R$^2$ 0.48--0.52 with negative coefficients on rainfall and humidity,
consistent with thermal stagnation and wet scavenging. The aggregate AQI sits at 39--51
(``moderate'') while the health-risk-adjusted index exceeds 90 seasonally --- \textbf{a
2-fold divergence between the two indices at identical concentrations}, which is the
paper's real contribution and an argument for HAQI-style reporting on any network serving
a PM-dominated airshed. Single station for a metropolitan area, and the provenance and
calibration history of a 12-year record in a data-sparse setting are not documented.}
{10.1371/journal.pgph.0006683}

\paperentry{DeMarsh et al.\ 2026 (backfill)}
{Community PurpleAir network versus regulatory monitors for PM$_{2.5}$ in Fresno County}
{Atmosphere}
{The SJV-CAIR network at UC Merced, built with community-based organisations, compared
against regulatory monitoring across Fresno County for June--July 2023. Community monitors
returned \textbf{concentration estimates comparable to regulatory instruments} where the
two co-locate, but identified \textbf{four weeks in which some locations episodically
exceeded the NAAQS that the regulatory network did not flag}, and interpolation maps built
with the denser network resolved spatial structure the regulatory-only map did not. Two
summer months in one county is a short and seasonally atypical window for the San Joaquin
Valley, and no formal co-location correction is reported. Carried as backfill --- Crossref
\texttt{created} 11 February 2026, outside this entry window --- and must be excluded from
weekly trend lines.}
{10.3390/atmos17020187}

\paperentry{Yukhymchuk et al.\ 2026}
{Low-cost particulate matter sensor performance under polluted conditions in Northeast China}
{Aerosol Air Qual Res}
{Crossref deposit of 7 August with \textbf{no abstract in Crossref, Europe PMC or Semantic
Scholar}. Carried on title, journal and DOI only; \textbf{no summary is offered because
none can be written without inventing content}. On title alone this is the single most
directly relevant record of the window --- LCS performance characterisation at the high
concentrations where light-scattering response is least well constrained --- and it should
be revisited as soon as an abstract appears.}
{10.1007/s44408-026-00159-6}

\paperentry{Sanders et al.\ 2026}
{Nonionic surfactant composition and hygroscopic growth of supermicron NaCl particles}
{Aerosol Sci Technol}
{Crossref deposit of 7 August, \textbf{metadata only --- no abstract deposited anywhere in
the retry chain}, so no summary is offered. Recorded here because surfactant control of
hygroscopic growth is the physical mechanism underlying humidity correction in
light-scattering PM sensors, and the supermicron size range is where nephelometric
response is most sensitive to growth factor.}
{10.1080/02786826.2026.2703042}

\paperentry{Huang et al.\ 2026}
{Metal ions and pH as controls on the optical properties of atmospheric brown carbon}
{Atmos Pollut Res}
{Crossref deposit of 7 August, \textbf{metadata only --- no abstract retrievable}, so no
summary is offered. Logged because brown-carbon absorption is the aerosol optical property
that most directly confounds the fixed refractive-index assumptions built into low-cost
optical particle counters.}
{10.1016/j.apr.2026.103167}

\band{Exposure assessment \& modelling\hfill 7 records}{Teal}

\paperentry{Myat et al.\ 2026}
{Ground, satellite and model assessment of PM$_{2.5}$ and attributable mortality in Myanmar}
{Environ Pollut}
{Multi-source assessment for a least-developed country with minimal regulatory monitoring:
two Yangon ground stations 2019--2024, TROPOMI HCHO and NO$_2$ nationally, and a
generalised additive model for meteorological dependence. Mean PM$_{2.5}$ was
\textbf{24.4\,$\mu$g/m$^3$ (SD 18.6)}, rising to \textbf{39.4\,$\mu$g/m$^3$} in the
November--February winter, with sulfate the leading measured component year-round and
nitrate and Ca$^{2+}$ seasonally. Compliance with the 5\,$\mu$g/m$^3$ WHO guideline would
have averted \textbf{2{,}106 (95\% CI 345--3{,}818) premature deaths annually} among
adults $\geq$30 in Yangon over 2020--2021. TROPOMI HCHO/NO$_2$ ratios indicate NO$_x$-
sensitive ozone chemistry nationally. \textbf{Two stations cannot represent a city of five
million}, and the health impact assessment inherits every assumption of the imported
concentration--response function.}
{10.1016/j.envpol.2026.128918}

\paperentry{Cai et al.\ 2026 (ACP)}
{Molecular characteristics and formation pathways of particulate organosulfur compounds in Shanghai}
{Atmos Chem Phys}
{Paired urban and suburban Shanghai sampling with high-resolution mass spectrometry,
resolving \textbf{1{,}964, 1{,}914 and 2{,}689 organosulfur molecular formulas} across the
sampled conditions and separating spatial, seasonal and day--night variation in formation
pathway. Organosulfates are ubiquitous but their production regime is not: the comparison
across contrasting sites is what isolates the pathways rather than the inventory itself.
The limitation that matters for downstream use is that molecular formula assignment from
exact mass does not fix structure, so pathway attribution is inferential. Relevant to
sensing as a concrete demonstration that a large, chemically distinct fraction of urban
PM$_{2.5}$ mass is invisible to any mass-only instrument.}
{10.5194/acp-26-11067-2026}

\paperentry{Blanco-Villafuerte et al.\ 2026 (backfill)}
{Hybrid low-cost-sensor, satellite and chemical-transport modelling of ambient PM$_{2.5}$ in Peru}
{Curr Environ Health Rep}
{Review of national PM$_{2.5}$ exposure-model construction in LMICs, anchored on the
NIH-funded Peru GeoHealth Hub's \textbf{176-monitor PurpleAir network} --- 62.5\% urban,
37.5\% rural, across all 24 regions, established with national environmental agencies. The
hybrid model fuses LCS, satellite retrievals, chemical transport output and machine
learning; initial application to Lima reached \textbf{R$^2$ = 0.88} against regulatory
monitors for 2010--2023, with extension to daily national 5\,km$^2$ fields for 2024--2026
in progress. R$^2$ against regulatory monitors in a single city is a weak validation for a
national product, and the review is written by the network's own investigators. The
deployment-ratio and institutional-integration detail is the transferable content.
Crossref \texttt{created} 5 January 2026 --- backfill, exclude from trend lines.}
{10.1007/s40572-025-00508-4}

\paperentry{Zheng et al.\ 2026}
{Spatially heterogeneous drivers of winter PM$_{2.5}$ in China from synchronous observation and explainable ML}
{Atmos Environ}
{Crossref deposit of 7 August, \textbf{metadata only --- Elsevier deposits no abstract to
Crossref and the record is not yet in Europe PMC, PubMed or Semantic Scholar}. No summary
is offered. Logged because synchronous multi-site observation paired with explainable
machine learning is the design that the sensing literature most often lacks: it is the
combination that could distinguish genuine spatial heterogeneity in drivers from
instrument heterogeneity across a network.}
{10.1016/j.atmosenv.2026.122285}

\paperentry{Rai et al.\ 2026}
{Land surface temperature, land cover and aerosol optical depth over Ahmedabad, 1993--2023}
{Environ Monit Assess}
{Multi-sensor remote sensing: Landsat LST and urban thermal field variance index for 1993,
2002, 2014 and 2023, with MODIS AOD and LST from 2002--2023 for the diurnal response.
\textbf{Urban extent grew by 171.15\,km$^2$ over 30 years} with sharp losses of fallow and
vegetated land, and the thermal environment polarised --- fallow land carries the worst
thermal comfort, water and vegetation the best, and NDVI/NDBI track the expected cooling
and warming signs. \textbf{The AOD--LST correlation is dual-signed}, which the paper
reports but does not mechanistically decompose into aerosol direct radiative effect versus
co-varying surface change; that confounding is the main threat to the inference. Read as
a heat-planning input rather than a PM exposure product --- AOD here is a columnar
quantity never converted to surface concentration.}
{10.1007/s10661-026-15748-6}

\paperentry{Meena et al.\ 2026}
{Transplanted lichen trace-metal accumulation along a Jaipur urban gradient}
{Environ Pollut}
{\textit{Phaeophyscia hispidula} thalli collected from a wildlife-sanctuary reference site
and transplanted to ten urban and one reference location in Jaipur for 90 days over the
March--June 2024 pre-monsoon, with Zn, Pb, Cd, Cr, Cu and Ni by ICP-MS and chlorophyll and
carotenoid concentrations by spectrophotometry. Urban sites accumulated more metal, and
\textbf{accumulation tracked chlorophyll loss while carotenoids rose} at several exposed
sites, consistent with a photoprotective adjustment. PCA and hierarchical clustering
resolved a spatial gradient. \textbf{No co-located instrumental PM measurement means the
biomonitor is never calibrated against a concentration}, so this yields a relative ranking
of sites, not an exposure estimate. Useful as an ultra-low-cost spatial screening layer
where no network exists; not a substitute for one.}
{10.1016/j.envpol.2026.128919}

\paperentry{Temu et al.\ 2026}
{A six-domain framework for lichen biomonitoring including symbiont and microbiome response}
{Environ Sci Pollut Res}
{Narrative synthesis proposing that lichen biomonitoring be extended from the classical
triad of pollutant accumulation, physiological injury and community composition to six
domains, adding functional traits, symbiotic interaction and lichen-associated microbiome
dynamics as \textbf{earlier-responding indicators} than bulk accumulation. The argument
that microbiome shifts precede visible physiological injury is plausible and is the part
worth testing, but the review supplies \textbf{no quantitative comparison of detection
limits or response latency across the six domains}, so the central claim of earlier warning
is asserted rather than demonstrated. The authors name geographic bias, species-specific
response and absent methodological standardisation as open problems --- the same
standardisation gap that limits low-cost sensor networks, and for structurally similar
reasons.}
{10.1007/s11356-026-38123-x}

\band{Cardiovascular \& metabolic\hfill 2 records}{Clay}

\paperentry{Li et al.\ 2026 (Environ Int)}
{PM$_{2.5}$ constituents, hepatic fat and stiffness in youth, modified by PNPLA3 genotype}
{Environ Int}
{113 Los Angeles Latino adolescents with obesity, MRI-measured hepatic fat fraction and
liver stiffness, serum metabolomics and lipidomics, and visit-year mean exposure to
\textbf{15 PM$_{2.5}$ components} at residence; linear regression plus g-computation for
mixture effects with PNPLA3 genotype as modifier, and moderated mediation. A simultaneous
IQR increase across all components was associated with \textbf{10.0\% higher hepatic fat
fraction}, with elemental carbon, bromine, copper, potassium and lead as leading
contributors --- a combustion and traffic signature. \textbf{Effects concentrated in PNPLA3
GG carriers}, who showed 550 mediating metabolite features against 196, broader pathway
enrichment, and the only significant total mediation. Replication in 81 mixed-ethnicity
young adults held for EC and Br. Small discovery sample, cross-sectional, modelled rather
than personal exposure --- the genotype interaction is the hypothesis, not the estimate.}
{10.1016/j.envint.2026.110438}

\paperentry{Zhao et al.\ 2026}
{Ischaemic stroke burden and attributable risk factors in adults aged 20--54, 1990--2021}
{Medicine (Baltimore)}
{GBD 2021 extraction for ischaemic stroke in the 20--54 age band, with age-period-cohort
models and population-attributable fractions by risk factor. In 2021 there were
\textbf{16{,}301{,}644 prevalent and 1{,}324{,}743 incident cases} (432.46 and 35.14 per
100{,}000), 105{,}859 deaths and 6.97 million DALYs. \textbf{Prevalence and incidence rose
over the period while mortality and DALYs fell}, with burden concentrated in high-middle
SDI regions, Eastern Europe, Central Asia and parts of North Africa and the Middle East.
Air pollution enters only as one of the GBD attributable risk factors, so this contributes
no exposure--response estimate and inherits every GBD modelling assumption, including
sparse primary data in the highest-burden regions. Carried for the age-band signal: young-
adult stroke burden is rising in exactly the regions where PM exposure is highest and
monitoring is thinnest.}
{10.1097/MD.0000000000050104}

\band{Burden, policy \& mitigation\hfill 3 records}{Violet}

\paperentry{Hou et al.\ 2026}
{Bimodal polysulfonamide fibrous aerogel for high-temperature PM$_{0.3}$ filtration}
{J Hazard Mater}
{Co-electrospun aerogel combining thermally stable PSA/PAN/SiO$_2$ coarse fibres with pure
PSA fine fibres in a 3D curly network at \textbf{$>$98\% porosity}, formed via
moisture-induced phase separation and jet whipping from the mechanical mismatch between
components. \textbf{PM$_{0.3}$ filtration efficiency 98.5\% at $\sim$84\,Pa pressure drop},
retained at \textbf{92.31\% and 73\,Pa after 12 h at 250\,$^\circ$C}, with acid and alkali
resistance and useful thermal insulation. The gap in the report is loading behaviour: no
long-duration dust-holding capacity or efficiency-versus-loading curve is given, and that
is what determines service life in a real industrial stack. Relevant to source-side
mitigation rather than monitoring, and to the ultrafine fraction that mass-based networks
under-weight.}
{10.1016/j.jhazmat.2026.143218}

\paperentry{Han et al.\ 2026}
{Air quality index product use versus clinician discussion among US children with asthma}
{AJPM Focus}
{Questionnaire administered inside the Transcript/Omics for Precision Medicine in Asthma
study, youth aged 6--21 recruited 2023--2025, with bivariate testing. \textbf{75\% of
caregivers were aware of an AQI or alert, 64\% checked it before going outdoors and 73\%
reported a behaviour change} --- but the delivery channel is consumer, not clinical or
public: a phone weather app in 75\% of cases, broadcast or social alerts 44\%, websites
24\%, county services only \textbf{14\%}. Only 41\% had ever discussed air pollution with a
clinician and 35\% received poor-air-quality recommendations. Self-report with no objective
behaviour measure and a single research cohort limit generalisability. The finding that
matters for network design: \textbf{the operative last mile is a commercial weather app,
and public monitoring agencies reach one family in seven}.}
{10.1016/j.focus.2026.100504}

\paperentry{Caikai et al.\ 2026}
{Risk-factor-attributable burden of respiratory infection and tuberculosis, 1990--2021}
{J Infect Dev Ctries}
{GBD 2021 across 204 countries, decomposing deaths, DALYs and age-standardised rates for
respiratory infection and TB across nine risk factors by age, sex and SDI, with APC and
AAPC trend estimation. Age-standardised mortality attributable to the traditional risks
fell steadily --- \textbf{air pollution $-$3.18\% annually}, malnutrition $-$4.46\%, alcohol
$-$3.35\% --- while metabolic-risk-attributable deaths rose in absolute terms despite
falling rates. Tobacco-attributable mortality was \textbf{3.6-fold higher in men}, and
children under 5 remain the group most exposed to environmental risks. Air pollution is a
GBD attributable fraction here, not a measured exposure--response, so this is a trend
context record: the attributable burden is falling on rate, which the accompanying
absolute counts do not support as a success narrative.}
{10.3855/jidc.21667}

\band{Other clinical endpoints\hfill 1 record}{Sage}

\paperentry{Sassano et al.\ 2026}
{Household solid-fuel use and kidney-disease mortality in the Golestan Cohort}
{Environ Health Perspect}
{50{,}045 adults aged 40--75 enrolled 2004--2008 in northeastern Iran and followed to April
2023, with validated fuel-use questionnaires and Cox models; 262 kidney deaths over
\textbf{724{,}064 person-years}. Per 10 additional years of household biomass use for
cooking or heating, \textbf{HR 1.20 (95\% CI 1.04--1.37)} for kidney-disease mortality;
kerosene 1.09 (0.95--1.24) and gas 1.00 (0.86--1.16) were null. \textbf{The biomass effect
splits on ventilation} --- 1.19 (1.06--1.34) without a chimney against 1.06 (0.95--1.18)
with one, $P_{\text{diff}}=0.025$ --- which is the strongest evidence in the paper that the
exposure is the smoke rather than a correlate of poverty, though confounding by
socioeconomic position is not fully excluded. \textbf{No PM measurement of any kind}: fuel
type and duration are the entire exposure model, and the ventilation interaction is
precisely the variable no ambient network can observe.}
{10.1021/EHP.6c00520}

\band{Mechanistic toxicology\hfill 1 record}{Deep}

\paperentry{Wu et al.\ 2026 (JHM)}
{TNF-$\alpha$ and EGF mediate endothelial activation downstream of PM$_{2.5}$ alveolar injury}
{J Hazard Mater}
{3D alveolar--vascular barrier constructed from BEAS-2B epithelium, THP-1 macrophages and
EA.hy926 endothelium, exposed to \textbf{100\,$\mu$g/mL SRM2786 PM$_{2.5}$ for 24 h}.
PM$_{2.5}$ entered epithelial cells by endocytosis, produced mitochondrial damage and
apoptosis, and raised TNF-$\alpha$, EGF, IL-8, TGF-$\alpha$ and IL-1$\beta$; extracellular
\textbf{TNF-$\alpha$ and EGF crossed the barrier and activated endothelial PI3K--Akt}, with
Nrf2 and CYP1A1 induction marking progressive oxidative stress. The model supplies a
mechanistic route from local alveolar injury to systemic endothelial dysfunction. The
single exposure concentration is far above any ambient equivalent and there is no
dose--response, so the result establishes the pathway's existence, not its relevance at
realistic exposures --- the standard and unresolved weakness of this literature.}
{10.1016/j.jhazmat.2026.143204}

\band{Reproductive \& developmental\hfill 1 record}{Sky}

\paperentry{GBD 2023 India Stillbirth Collaborators 2026}
{State-level stillbirth rates in India at $\geq$22 and $\geq$28 weeks, GBD 2023}
{Lancet Reg Health Southeast Asia}
{First GBD estimation of stillbirth rate under the WHO ICD-11 $\geq$22-week definition
alongside the conventional $\geq$28-week cutoff, disaggregated to 31 Indian geographic
units. \textbf{122 data sources and 2{,}684 location-year observations} were modelled by
spatiotemporal Gaussian process regression on the stillbirth-to-neonatal-mortality ratio,
anchored on GBD 2023 fertility and all-cause neonatal mortality, with 95\% uncertainty
intervals and a secondary analysis against SDG3 indicators. \textbf{India carries the
largest global stillbirth burden at both cutoffs}, and the $\geq$28-week convention
materially understates it. Air pollution enters only through GBD's risk-factor
architecture, so no PM exposure--response is estimated here. Carried because state-level
disaggregation is the resolution at which a PM--stillbirth analysis could be attempted, and
because it defines the denominator such an analysis would use.}
{10.1016/j.lansea.2026.100819}

\clearpage

% ==================================================================== TABLE
\section*{\sffamily\bfseries\color{Deep}Corpus register}
\vspace{-1.5mm}{\color{Amber}\rule{\linewidth}{1.1pt}}\vspace{2.5mm}

{\fontsize{6.7}{8.4}\selectfont
\setlength{\tabcolsep}{2.4pt}
\begin{longtable}{@{}>{\RaggedRight}p{25mm}>{\RaggedRight}p{41mm}>{\RaggedRight}p{27mm}%
>{\RaggedRight}p{23mm}>{\RaggedRight}p{18mm}>{\RaggedRight}p{15mm}@{}}
\toprule
\textsf{\textbf{First author}} & \textsf{\textbf{Focus}} & \textsf{\textbf{Journal}} &
\textsf{\textbf{Design}} & \textsf{\textbf{Metric}} & \textsf{\textbf{Region}}\\
\midrule
\endfirsthead
\toprule
\textsf{\textbf{First author}} & \textsf{\textbf{Focus}} & \textsf{\textbf{Journal}} &
\textsf{\textbf{Design}} & \textsf{\textbf{Metric}} & \textsf{\textbf{Region}}\\
\midrule
\endhead
\multicolumn{6}{@{}l@{}}{\cellcolor{Mist}\textsf{\textbf{\textcolor{Deep}{Neuro / mental health}}}}\\[0.6mm]
\href{https://doi.org/10.1001/jamanetworkopen.2026.25015}{Dales} & Wildfire-sourced PM$_{2.5}$ and ED visits for migraine and other primary headache & \textit{JAMA Netw Open} & Case-crossover & PM$_{2.5}$ (wildfire) & Canada\\
\href{https://doi.org/10.1093/ageing/afag215}{Glover} & Umbrella meta-review of urban environmental exposures and cognitive health & \textit{Age Ageing} & Umbrella review / meta-analysis & PM$_{2.5}$ & Global\\
\href{https://doi.org/10.1016/j.nsa.2026.107019}{Goossens} & In utero air pollution exposure and pre-/postnatal brain development & \textit{Neurosci Appl} & Narrative review & PM$_{2.5}$ / PM$_{10}$ / PM$_{\mathrm{coarse}}$ & Global\\
\href{https://doi.org/10.1016/j.isci.2026.116617}{Ma} & PM$_{2.5}$-induced neurotoxicity and depression: exposure risks and mechanisms & \textit{iScience} & Narrative review & PM$_{2.5}$ + components & Global\\
\href{https://doi.org/10.3760/cma.j.cn112150-20250804-00752}{Pan} & Maternal PM exposure and offspring neurodevelopment: effects and mechanisms & \textit{Zhonghua Yu Fang Yi Xue Za Zhi} & Narrative review & PM (bulk) & China\\
\addlinespace[0.9mm]
\multicolumn{6}{@{}l@{}}{\cellcolor{Mist}\textsf{\textbf{\textcolor{Teal}{Cardiovascular \& metabolic}}}}\\[0.6mm]
\href{https://doi.org/10.3760/cma.j.cn112338-20251120-00843}{Huang} & PM$_{2.5}$ and dyslipidaemia in rural Chinese children; dietary effect modification & \textit{Zhonghua Liu Xing Bing Xue Za Zhi} & Cross-sectional (multistage) & PM$_{2.5}$ & China\\
\href{https://doi.org/10.32394/pe/218966}{Kazmierska} & Air pollution and cardiovascular disease: review of current data & \textit{Przegl Epidemiol} & Narrative review & PM$_{2.5}$ / PM$_{10}$ & Global\\
\href{https://doi.org/10.1007/s00125-026-06814-2}{Xi} & Air pollution, built environment and temperature vs. prediabetes phenotypes & \textit{Diabetologia} & Prospective cohort & PM$_{2.5}$ & Germany\\
\href{https://doi.org/10.1161/JAHA.125.046917}{Yang} & Short-term PM$_{2.5}$ and hospitalisation/ED use in atrial fibrillation & \textit{J Am Heart Assoc} & Case-crossover & PM$_{2.5}$ (wildfire-influenced) & USA\\
\addlinespace[0.9mm]
\multicolumn{6}{@{}l@{}}{\cellcolor{Mist}\textsf{\textbf{\textcolor{Sky}{Reproductive \& developmental}}}}\\[0.6mm]
\href{https://doi.org/10.1016/j.jes.2025.10.043}{Chen} & Prenatal air pollution, fetal growth trajectory phenotypes and child neurodevelopment & \textit{J Environ Sci (China)} & Prospective cohort & PM$_{2.5}$ / PM$_{10}$ & China\\
\href{https://doi.org/10.1016/j.ecoenv.2026.120566}{Huang} & Joint prenatal bisphenol + air-pollutant mixture exposure and preschool IQ & \textit{Ecotoxicol Environ Saf} & Prospective cohort (mixtures) & PM$_{2.5}$ / PM$_{10}$ & China\\
\href{https://doi.org/10.1007/s40572-026-00552-8}{Peterson} & Ambient air pollution, PAHs and pubertal development: critical review & \textit{Curr Environ Health Rep} & Critical review & PM$_{2.5}$ components & Global\\
\href{https://doi.org/10.1016/j.envres.2026.125313}{Xie} & Long-term PM$_{2.5}$ and timing of pubertal development in boys and girls & \textit{Environ Res} & Prospective cohort & PM$_{2.5}$ & China\\
\addlinespace[0.9mm]
\multicolumn{6}{@{}l@{}}{\cellcolor{Mist}\textsf{\textbf{\textcolor{Sage}{Respiratory \& allergic}}}}\\[0.6mm]
\href{https://doi.org/10.1007/s00484-026-03289-z}{Bayar Muluk} & Climate change and otolaryngology: airway disease impacts and emission reduction & \textit{Int J Biometeorol} & Narrative review & PM$_{2.5}$ & Global\\
\href{https://doi.org/10.4103/mgr.MEDGASRES-D-25-00164}{Liao} & Short-term air pollution, respiratory symptoms and lung function in asthma & \textit{Med Gas Res} & Panel study & PM$_{2.5}$ / PM$_{10}$ & China\\
\href{https://doi.org/10.3390/epidemiologia7040097}{Sanchez de Prada} & Ambient PM and asthma-coded admissions in children 0-3 y, Spain & \textit{Epidemiologia} & Ecological panel & PM$_{2.5}$ / PM$_{10}$ & Spain\\
\href{https://doi.org/10.1016/j.jhazmat.2026.143076}{Yang} & Allergenic vs. toxic risk of airborne seed hairs compared with ambient PM$_{2.5}$ & \textit{J Hazard Mater} & Field measurement & PM$_{2.5}$ + bioaerosol & China\\
\addlinespace[0.9mm]
\multicolumn{6}{@{}l@{}}{\cellcolor{Mist}\textsf{\textbf{\textcolor{Violet}{Mechanistic toxicology}}}}\\[0.6mm]
\href{https://doi.org/10.1016/j.jes.2025.12.074}{Ding} & PM$_{2.5}$ promotes vascular calcification via miR-27b-5p/HES1 signalling & \textit{J Environ Sci (China)} & In vitro + organoid + murine & PM$_{2.5}$ & China\\
\href{https://doi.org/10.1007/s00360-026-01698-5}{Maraschi} & Osmoregulatory disruption and metal bioaccumulation from settleable atmospheric PM & \textit{J Comp Physiol B} & Controlled exposure (ecotox) & Settleable APM & Brazil\\
\href{https://doi.org/10.1016/j.jes.2025.11.018}{Wu} & Macrophage NLRP3 inflammasome in real-ambient PM-induced hepatic glucose dysregulation & \textit{J Environ Sci (China)} & Real-ambient murine + co-culture & Real-ambient PM & China\\
\addlinespace[0.9mm]
\multicolumn{6}{@{}l@{}}{\cellcolor{Mist}\textsf{\textbf{\textcolor{Clay}{Occupational \& indoor}}}}\\[0.6mm]
\href{https://doi.org/10.3389/fmedt.2026.1796828}{Chen} & Nested suction instrument to cut intraoperative surgical-smoke PM$_{2.5}$ & \textit{Front Med Technol} & Bench / device experiment & PM$_{2.5}$ (surgical smoke) & China\\
\href{https://doi.org/10.1093/annweh/wxag055}{Cole} & Determinants of respirable crystalline silica exposure in Australian tunnelling & \textit{Ann Work Expo Health} & Observational exposure study & Respirable dust / RCS & Australia\\
\href{https://doi.org/10.1016/j.jhazmat.2026.143081}{Ding} & Stage-resolved hair Zn-isotope dynamics as a firefighting combustion-exposure tracer & \textit{J Hazard Mater} & Repeated-measures biomarker & Combustion-derived PM & China\\
\href{https://doi.org/10.3390/jox16040119}{Ramsden} & Xenobiotic hazards in aircraft cabin air & \textit{J Xenobiot} & Review + risk economics & Ultrafine particles & Global\\
\addlinespace[0.9mm]
\multicolumn{6}{@{}l@{}}{\cellcolor{Mist}\textsf{\textbf{\textcolor{Amber}{Sensing, forecasting \& instrumentation}}}}\\[0.6mm]
\href{https://doi.org/10.1016/j.jenvman.2026.130515}{Shi} & Adaptive spatiotemporal graph transformer for multi-site multi-step PM$_{2.5}$ forecasting & \textit{J Environ Manage} & Deep-learning model & PM$_{2.5}$ & China / India\\
\href{https://doi.org/10.3390/bios16070380}{Tsuchiya} & Seasonal and indoor-outdoor characteristics of fluorescent aerosol particles in offices & \textit{Biosensors} & Real-time sensor campaign & Bioaerosol / coarse & Japan\\
\addlinespace[0.9mm]
\multicolumn{6}{@{}l@{}}{\cellcolor{Mist}\textsf{\textbf{\textcolor{Amber}{Exposure assessment \& modelling}}}}\\[0.6mm]
\href{https://doi.org/10.1029/2025GH001621}{Solorzano-Araque} & Multiscale exposure modelling: LOTOS-EUROS + mobility surveys + particle toxicity & \textit{GeoHealth} & CTM + mobility integration & PM$_{2.5}$ / PM$_{10}$ & Colombia\\
\addlinespace[0.9mm]
\multicolumn{6}{@{}l@{}}{\cellcolor{Mist}\textsf{\textbf{\textcolor{Coral}{Burden, policy \& mitigation}}}}\\[0.6mm]
\href{https://doi.org/10.1002/hsr2.72886}{Karamzad} & Ischaemic heart disease burden in MENA 1990-2021; PM attributable fraction 28.2\% & \textit{Health Sci Rep} & GBD secondary analysis & Ambient PM & MENA\\
\href{https://doi.org/10.1371/journal.pone.0347989}{Rees} & HEPA filtration to reduce respiratory infection in care homes & \textit{PLoS One} & RCT process evaluation & Filtered PM & UK\\
\href{https://doi.org/10.1016/j.jes.2026.01.009}{Shan} & Policy-driven emission reductions and post-pandemic rebound, high-resolution inventories & \textit{J Environ Sci (China)} & Emission inventory & PM emissions & China\\
\addlinespace[0.9mm]
\multicolumn{6}{@{}l@{}}{\cellcolor{Mist}\textsf{\textbf{\textcolor{Slate}{Other clinical endpoints}}}}\\[0.6mm]
\href{https://doi.org/10.3390/dj14070440}{Lupita} & PM$_{2.5}$, preventive-care capacity and caries/edentulism burden across the EU & \textit{Dent J} & Ecological panel & PM$_{2.5}$ & EU-27\\
\href{https://doi.org/10.1016/j.jdin.2026.06.008}{Park} & Long-term PM exposure and risk of androgenetic alopecia & \textit{JAAD Int} & Nationwide cohort & PM$_{2.5}$ / PM$_{10}$ & South Korea\\
\href{https://doi.org/10.3171/2026.3.PEDS2615}{Sun} & Environmental exposure and area deprivation vs. paediatric hydrocephalus outcomes & \textit{J Neurosurg Pediatr} & Registry cohort + geospatial & PM$_{2.5}$ (satellite-derived) & USA\\
\addlinespace[0.9mm]
\bottomrule
\end{longtable}}


\vspace{4mm}
\begin{keybox}
{\sffamily\bfseries\footnotesize\color{Deep}Provenance and method}\\[1.2mm]
{\fontsize{7.2}{8.8}\selectfont
Entry window \textbf{7 August 2026}, contiguous with the 6 August issue; no gap. Every
PubMed record shipped here was verified by \texttt{efetch} to carry a
\texttt{PubMedPubDate[@PubStatus="pubmed"]} of \textbf{7 August 2026} --- 27 of 27, so no
record was held for the next issue and none was re-indexed from an earlier window. Two
Consensus records are dated January and February 2026 by Crossref \texttt{created} and are
\textbf{explicitly flagged backfill} in their entries; they must be excluded from any
weekly trend line.
\textbf{Sources.} PubMed on two separate axes windowed on \texttt{[EDAT]} --- the connector
returned 19 health and 17 instrumentation hits (union 27), and \textbf{the local harvester
returned 18 health and 0 sensing}, its sensing query producing an empty set for the second
time. The connector remains the better recall leg on the instrumentation axis and both
were merged before screening. Union 27, 1 dropped as an already-shipped duplicate, 18
carried, 8 rejected out of scope (Arctic marine particulate organic matter, indoor VOC
fabric-air partitioning, cigarette-smoke bacteriology, a Legionella water assay, cannabis
analytical chemistry, heat-not-burn fire-safety testing and two editorials with no primary
data). Europe PMC on \texttt{CREATION\_DATE} returned 1, already in the PubMed set.
\textbf{Crossref scoped by journal ISSN across eight aerosol and atmospheric titles PubMed
does not index returned 10 --- the most productive sensing leg again} --- of which 6 were
carried and 4 rejected out of scope (GNSS refractivity, frost-point hygrometry, thunderstorm
polarimetric radar, and a pollen/allergic-rhinitis study excluded as aeroallergen rather
than mass-based PM). Consensus on three axes (calibration and drift; PurpleAir/Plantower
network siting; satellite--surface fusion) returned 9 hits capped at 3 per sweep by the
free tier; \textbf{7 resolved to DOIs already in \texttt{seen.json}} and 2 were carried as
backfill. arXiv returned 60 records, \textbf{none published on or after 1 August}, so the
recency screen carried nothing. OpenAlex returned \texttt{HTTP 429} on top of the standing
paid-plan block --- that leg has now failed on every run since 28 July.
\textbf{Screening.} Deduplication keys PMID $\rightarrow$ lowercased DOI $\rightarrow$
title signature, checked against \texttt{seen.json} \emph{before} summarising; Consensus
titles were resolved to DOIs through Crossref before screening, which is what caught the
seven duplicates. \textbf{Four records are carried on metadata alone} --- DOIs verify in
Crossref but no abstract exists in Crossref, Europe PMC or Semantic Scholar --- and are
marked in their entries; \textbf{no summary was written from a title}. First-author
surnames and DOIs were read from \texttt{efetch}, taking
\texttt{ELocationID[@EIdType='doi']} before \texttt{ArticleIdList} so that no reference-list
DOI is picked up.
\textbf{Labels.} Subtopic, design, particle-metric and geography assignments are
single-label judgements from title, abstract, keywords and MeSH terms, not a controlled
vocabulary. The sixteen effect estimates are transcribed verbatim; one is a co-pollutant
(NO$_2$) rather than PM and is marked as such in the forest caption and the register.
Jiang et al.'s PM$_{10}$ length-of-stay RR of 1.63 (0.18--15.94) is deliberately excluded
from the forest plot: its interval spans an order of magnitude and would compress every
other estimate to a point.
All DOI links resolve to the publisher of record and passed the automated DOI gate before
this issue compiled. Abstracts remain the copyright of their publishers.\par}
\end{keybox}

\end{document}
